\chapter{Theoretical Background and Related Work\label{cha:chapter2}}

This chapter provides the background knowledge and literature context required to understand the prompt-sensitivity concepts and empirical methods applied in this thesis.

\section{Prompt Sensitivity and Prompt Engineering in LLMs\label{sec:prompt-sensitivity}}

Since large language models with strong few-shot and instruction-following abilities appeared in 2020, prompt engineering has become an essential part of working with these models \cite{brown2020fewshot}. Recently, prompting techniques have evolved from simple task descriptions to include instructions, examples, input/output format, role descriptions, reasoning cues, and conversational style. Recent surveys describe prompt engineering not as a collection of standard rules but as a broad and heterogeneous design space with many proposed techniques \cite{schulhoff2024promptreport, zhao2023llmsurvey}. Systematic reviews in educational contexts further examine how prompt engineering practices are being integrated into curricula \cite{lee2025promptedu}. A growing line of automatic prompt-optimization work instead treats prompts as searchable or compilable artifacts \cite{pryzant2023apo, yang2023opro, khattab2023dspy}.

Early research on prompt-programming and few-shot calibration work also showed that prompts are not merely presentation layers but can significantly influence model behaviors. Reynolds and McDonell describe prompts as natural-language programs that communicate task intention and structure \cite{reynolds2021promptprogramming}. Zhao et al. demonstrated that prompt format, example choice, and example order can create a large variance in few-shot accuracy \cite{zhao2021calibrate}. Lu et al. found that prompt example permutations alone can change performance from near-random to near state-of-the-art \cite{lu2021fantastically}. This is consistent with findings on in-context learning as implicit, context-dependent reasoning, as shown in \cite{xie2022icl}. Subsequent work has extended these findings to superficial prompt features. Indeed, Sclar et al. found that minor changes in format can be responsible for significant differences in performance, such as spacing or delimiter choice \cite{sclar2023spurious}. Zhu et al. studied model robustness to adversarial perturbations of prompts \cite{zhu2023promptrobust}, while Zhuo et al. analyzed the effect of prompt formats on code-generation evaluation \cite{zhuo2024prosa}. Finally, even automatically discovered prompts that read as nonsense to humans can be better than manually created prompts, as reported in \cite{battle2024eccentric}. All these findings give reason to treat prompts as elements of the experimental design, rather than as incidental text.

Nevertheless, most of the studies in prompt-sensitivity examine primarily the aspect of correctness. While correctness is an important and necessary aspect in a production software application, it is insufficient. There can be correct answers that take too long or are overly complex for quick interpretation. This supports one of the ideas in behavioral testing of NLP systems, which is that other characteristics besides correctness are also important \cite{ribeiro2020checklist}. Accordingly, one of the points made in this thesis is that correctness has to be distinguished from the output behavior measured with the help of certain indicators such as token count, time of completion, result-only compliance, empty responses, and extraction failures.

The social context and tone are also taken into consideration in this area of study. Research on emotional stimuli suggests that emotional discourse may enhance the quality of generated messages \cite{li2023emotionprompt}, whereas research on politeness shows that it influences the form and content of the messages being generated \cite{yin2024politeness}. Another element from the studies on sycophancy is that the output of bots can reflect the praise, agreement, or pressure to comply that the user exhibits \cite{sharma2023sycophancy}. The phenomenon of role-playing has also been studied as a way to influence the style and content of communication \cite{shanahan2023roleplay}. Current research is focused on aligning the output of LLMs with the user’s needs through various techniques of prompting and fine-tuning \cite{shen2023alignment}. These findings suggest that tone should be taken into account in this area of study, but they do not provide general principles. In other words, a polite prompt can produce different outcomes in different contexts.

\section{LLMs for Code Generation and Software-Engineering Tasks}

Code generation is the main software engineering application of LLMs. Benchmarks such as HumanEval and MBPP evaluate whether generated Python programs satisfy hidden or visible tests \cite{chen2021codex, austin2021mbpp}. For more realistic scenarios, SWE-bench extends evaluation to the repository level \cite{jimenez2023swebench}, and subsequent benchmarks broaden the picture to diverse function calls and complex instructions \cite{zhuo2024bigcodebench} as well as to many programming languages \cite{cassano2023multipl}. Recent open-source code models such as CodeLlama, StarCoder 2, WizardCoder, and Magicoder have substantially advanced the state of the art in code generation capabilities \cite{roziere2023code, lozhkov2024star, luo2024wizard, wei2024magicoder}. Comprehensive surveys of code generation methods and challenges provide additional context for understanding the current state of the art \cite{jiang2024code}. Those evaluations measure function correctness objectively rather than subjective preference. They also reveal a practical problem in evaluation: they depend on the process of extracting and executing code from model responses; rigorous re-evaluations show that execution-based accuracy is often lower than headline numbers suggest \cite{liu2023nece}. Therefore, generated code must be well recognizable, syntactically correct, and compatible with the test suites.

For these reasons, code generation is particularly suitable for studying prompt framing. In programming workflows, useful products are often not explanations, but functions, patches, commands, or structured outputs. Prompt framing affects whether the model returns clean code, code wrapped with markdown, text with code, or malformed output. A prompt can therefore change the pass rate by changing parser compatibility, even when the model's ability to solve the task is unchanged.

This thesis uses small, controlled Python tasks rather than repository-level tasks: task difficulty stays constant while framing changes, so prompt effects can be isolated. Chapter~\ref{cha:chapter3} discusses what this design does and does not represent.

\section{Structured Prompts, Output Constraints, and Parser Compatibility}

In LLM applications, it is common to use structured prompts and output constraints. According to \cite{liu2024structuredoutput}, developers use JSON, XML, schemas, field names, or strict response instructions when framing prompts. Although structured output improves the compatibility of model outputs with software systems, it does not guarantee positive results. Merely placing inputs inside XML/JSON formats does not ensure that the model will produce valid structured output, and overly complicated formatting may result in the model losing helpful contextual information. A different technique suggests using a more complex form of grammar-based generation at the decoding stage \cite{geng2023grammar}, but this thesis intentionally excludes this method from the analysis to isolate the prompt-side framing approach.

In this work, plaintext, XML, and JSON formats are tested as prompt framing methods. The purpose of this analysis is not just to identify whether a structured format brings about positive results, but also whether it changes the behavior and correctness of the model, and whether such changes depend strongly on the model and task being used (Python code vs. math).

This difference is important because users often confuse input wrapping with output control. Users may provide their tasks in the form of JSON while asking the model to create code or final numbers. In this regard, wrapping is a kind of input structure rather than ensuring a certain output form. In the paper, this problem is dealt with by monitoring both the correctness (pass rate) and the structural compliance (e.g., result-only or code-only) of the answers.

In the field of software engineering, this difference influences the way engineers construct pipelines rather than simply their terminology. Using JSON/XML to wrap prompts improves readability but doesn’t guarantee the response format of the model. The model may still generate text, markdown, code, or incorrect structures. So, downstream applications require both input formatting restrictions and output validation processes. This thesis classifies the wrapping formats as a subcomponent of the entire model-tool interface.

\section{Evaluation of Reasoning and Final-Answer Extraction}

Benchmarks in mathematical reasoning such as GSM8K evaluate whether models can find correct answers to word problems \cite{cobbe2021gsm8k}, whereas MATH features competition-level questions \cite{hendrycks2021math}. Studies on reasoning techniques such as zero-shot reasoning and chain-of-thought have shown that prompts affect the flow of multi-step reasoning \cite{wei2022cot, kojima2022zeroshot}; this work has been extended in self-consistency sampling \cite{wang2022selfconsistency}, least-to-most decomposition \cite{zhou2022leasttomost}, and search-based reasoning \cite{yao2023tot}. Recent reinforcement learning models, such as DeepSeek-R1, demonstrate how the use of this technique can improve multi-step problem solving \cite{guo2025deepseek}. Process-supervision studies also show that evaluating intermediate steps rather than just final answers leads to better results in mathematics \cite{lightman2023verify, uesato2022process}. Recent studies concerning reasoning techniques have looked at what these models can achieve regarding problem complexity \cite{shojaee2025illusion}. However, it still takes considerable time to evaluate answers to problems, as it is necessary to perform extraction, normalization, and comparison. Furthermore, reasoning traces are unlikely to correspond to what the model has actually computed \cite{turpin2023unfaithful}.

In this thesis, math problems serve as a secondary validation domain with a different output contract: code tasks reward executable artifacts and parser compatibility, whereas math tasks require deriving a correct answer even when the final output must be compact.

The mathematical perspective helps to differentiate two concepts that are easily confused. Instructions requiring only the final answer may lead to easier evaluation; however, this does not imply that deriving the solution is simpler. Reasoning visible in the output does not indicate that the internal thought process behind it is faithful. Thus, this research stipulates that longer responses are not automatically considered better, and shorter ones are not automatically considered worse. The main conclusions show the relationships observed between the prompt format, the length of the response, the level of compliance, and the accuracy of the answer.

\section{Research Gap\label{sec:research-gap}}

Table~\ref{tab:related-work-comparison} illustrates how this master thesis relates to previous work within pertinent dimensions. Three principal observations can be made. First, existing studies in the area mainly consider isolated variables when examining the issue of prompt sensitivity. For instance, Sclar et al. \cite{sclar2023spurious} analyze the impact of surface formatting, Lu et al. \cite{lu2021fantastically} work with the order of examples, etc. Nevertheless, it is worth noting that no study has considered tone, role, and wrapper format together. Therefore, their possible interactions remain unexplored. Second, on the topic of prompt sensitivity, previous works have also explored the importance of code generation in studies \cite{zhuo2024prosa}; however, their findings have not been incorporated into the context of reasoning and assessment results other than accuracy, nor into mediation-style analyses of the observed differences. Third, broad multi-metric evaluations like HELM \cite{liang2023helm} measure behaviors beyond accuracy but do not isolate prompt framing as the experimental variable. Furthermore, the most relevant tone studies report only accuracy, leaving the cost, latency, and compliance aspects of prompt framing unquantified.

{\footnotesize
\setlength{\tabcolsep}{2.5pt}
\begin{longtable}[]{@{}lccccccccc@{}}
\caption{Comparison with the closest prior studies. \checkmark{} = included, $\sim$ = partial or single-family, -- = not included. Behav. = behavioral outcomes beyond accuracy (tokens, latency, compliance).}\label{tab:related-work-comparison}\tabularnewline
\toprule\noalign{}
Study & Multi & Tone & Role & Fmt & Code & Reas. & Behav. & Fact. & Repro. \\
\midrule\noalign{}
\endfirsthead
\toprule\noalign{}
Study & Multi & Tone & Role & Fmt & Code & Reas. & Behav. & Fact. & Repro. \\
\midrule\noalign{}
\endhead
\bottomrule\noalign{}
\endlastfoot
Zhao et al.~\cite{zhao2021calibrate} & $\sim$ & -- & -- & \checkmark & -- & -- & -- & -- & -- \\
Lu et al.~\cite{lu2021fantastically} & $\sim$ & -- & -- & $\sim$ & -- & -- & -- & -- & $\sim$ \\
Sclar et al.~\cite{sclar2023spurious} & \checkmark & -- & -- & \checkmark & -- & -- & -- & $\sim$ & \checkmark \\
Zhu et al.~\cite{zhu2023promptrobust} & \checkmark & -- & -- & $\sim$ & -- & -- & -- & -- & \checkmark \\
Li et al.~\cite{li2023emotionprompt} & \checkmark & $\sim$ & -- & -- & -- & $\sim$ & -- & -- & $\sim$ \\
Yin et al.~\cite{yin2024politeness} & \checkmark & \checkmark & -- & -- & -- & $\sim$ & -- & $\sim$ & $\sim$ \\
Zhuo et al.~\cite{zhuo2024prosa} & \checkmark & -- & -- & \checkmark & \checkmark & -- & -- & $\sim$ & $\sim$ \\
Liang et al.~\cite{liang2023helm} & \checkmark & -- & -- & -- & \checkmark & \checkmark & \checkmark & -- & \checkmark \\
\textbf{This thesis} & \checkmark & \checkmark & \checkmark & \checkmark & \checkmark & \checkmark & \checkmark & \checkmark & \checkmark \\
\end{longtable}
}

Beyond these studies, this thesis also contributes: (i) a comprehensive factorial assessment of tone, role framing, and format (18 variations), thus measuring each factor against a stable baseline instead of in isolation; (ii) an analysis of 8 models ranging from local Ollama deployment systems to hosted APIs, treated as complementary settings rather than pooled benchmarks; (iii) an evaluation of executable code as well as math verification for the final answer under the same prompt matrix; (iv) behavioral outcomes (tokens, time, result-only compliance, and failure modes) reported together with pass@1 and a mediation-style decomposition in associational terms; and (v) mixed-effects inferential analysis, dominant-model sensitivity checks, and public replication materials. The approach is useful in terms of understanding how prompt framing functions, but not as a general mechanism for ensuring reliability.
